\chapter{Future Work}
\label{ch:future-work}

% Promoted from section 10.3 on the first examiner's 2026-08-18 review
% (registry row R37): the section had accumulated a series of very short
% numbered subsections, which read better as sections of their own chapter.
% Labels are carried over unchanged so every existing reference resolves.

The items below are ordered by what the results of this thesis suggest, not by convenience.
The first follows directly from the finding that reading is the binding constraint.

\section{Improve the Reading Stage, Not the Structuring Stage}
\label{sec:fw-reader}

The attribution result (Section~\ref{sec:results-attribution}) says the largest available gain is in
reading. Three specific directions follow.

\textbf{Adapt the reading model rather than the structuring model.} This thesis adapted the stage
that was not the bottleneck. Low-rank adaptation of the vision tower on invoice pages is the
untested experiment that the attribution evidence actually recommends, and it is a different
experiment from the one reported here.

\textbf{Ensemble across readers with complementary failure modes.} The two finalists fail
differently and in a way that looks exploitable: one drops letterhead and footer regions while
reading characters near-flawlessly, the other reads every region with occasional character slips
(Section~\ref{sec:results-reader}). A combination that takes regions from the first and characters from
the second is a concrete proposal, and the per-miss classification data needed to evaluate it
already exists.

\textbf{Pre-process before reading.} Deskewing, contrast normalisation and resolution
enhancement address the channel gap directly. Given that the channel gap is the dominant
real-world cost measured here, an intervention on the image may be worth more than any
intervention on the models.

\section{Teach Abstention Rather Than Always Answering}
\label{sec:fw-abstention}

The adaptation failed by raising spurious emission (Section~\ref{sec:disc-finetune}), which is what
supervised fine-tuning on exclusively-filled targets would be expected to do
\parencite{uluoglakci2026humility}. The corresponding untested variant is a training set that includes
appropriate absence as a taught behaviour: documents where a field is genuinely missing, with the
correct answer being the null. This thesis's training data contained no such examples, and the
mechanism points directly at that gap.

\section{Knowledge-Graph Construction}
\label{sec:fw-kg}

The most direct continuation of the designed-but-unbuilt architecture
(Section~\ref{sec:design-layer2}). The interesting question is not whether a graph can be built --- it
can --- but how extraction error propagates into it: does a field-level error rate predict a
compliance-decision error rate, or does the graph's redundancy absorb some errors while amplifying
others? That is a measurement question, and the per-field outcome data this thesis produces is the
right input to it.

\section{Analytical Query and Retrieval Comparison}
\label{sec:fw-query}

Building the query layer (Section~\ref{sec:design-layer3}) would let the second and third unevaluated
hypotheses be tested. The open question worth answering is narrower than the design: which
\emph{classes} of question genuinely benefit from graph structure over flat retrieval, since the
answer is unlikely to be uniform across question types.

\section{A Cloud Comparison That Respects the Premise}
\label{sec:fw-cloud}

Quantifying the accuracy cost of the local-only constraint requires running the same corpus
through a frontier cloud model. That cannot be done with client documents without violating the
premise. It can be done on the synthetic corpus, where every document is generated and none is
confidential --- and since the synthetic corpus carries exact reference data, such a comparison
would be cleanly scored. The result would be a lower bound on the cost of the constraint, because
the synthetic corpus is the easy case for both sides.

\section{Extend the Answer Key: Real Line Items, Local Adjudication}
\label{sec:fw-local-gt}
\label{sec:fw-groups}

Two pieces of answer-key work would extend what the held-out corpus can support, and they
share the same adjudication discipline.

\textbf{Measure real line items.} The route to closing Section~\ref{sec:lim-groups} is known
and is not research: adjudicate the held-out group rows under a \emph{set-based} matching
rule instead of a positional one, so that a reader which merges two lines is not penalised
as though it had misread them. That requires a matching criterion defined before
adjudication begins, and it requires the same ranked-escalation and provenance-class
discipline the header key used. It is expensive but it is bounded work.

\textbf{Adjudicate locally.} Closing Section~\ref{sec:lim-cloud-gt} requires the three-channel
adjudication to run without a cloud
service: a local vision model as one channel, a local \ac{OCR}-plus-layout stack as a second, the
embedded text layer as a third, and the same printed-text gate and ranked escalation. Whether
three \emph{local} channels are independent enough for the two-channel-agreement rule to carry
its current warrant is itself an open question, and it is the one that would have to be answered
first.

\section{Measure the Floor of the Envelope}
\label{sec:fw-envelope}

The most immediately available missing number in this thesis, and the cheapest to obtain: read
the 39 held-out invoices with the selected reader \emph{on the laptop}, under the pixel cap the
floor's runtime imposes, then structure and score them with the pipeline and scoring rules used
throughout. Nothing new has to be built --- the reading script already supports the local path
and already documents the cap as the reason it exists --- and nothing leaves the premises, so
unlike the answer-key work (Section~\ref{sec:lim-cloud-gt}) this measurement is itself deployable
inside the constraint it evaluates.

The comparison against Table~\ref{tab:heldout-headline} would be clean, because only the
reading venue and resolution change: same documents, same answer key, same structurer, same
scorer, same greedy decoding. The result would convert Section~\ref{sec:lim-hardware} from a stated
gap into a measured quantity --- the accuracy cost of the floor relative to the target class,
on real documents --- which is the number a company equipping below the target class would
actually want and the one this thesis is currently unable to give. Two protocol requirements follow from
Chapter~\ref{ch:measurement-validity}: the local transcripts must be archived like every other
generation, and the comparison must be pre-registered as a two-cell design before the run, so
that a disappointing local figure cannot be re-framed afterwards.

\section{Robustness on Degraded Documents}
\label{sec:fw-robustness}

The channel gap measured here is the most actionable finding in the thesis, and it was measured on
ten photographed documents. A controlled degradation study --- the same documents captured at
several resolutions, angles and lighting conditions --- would turn a two-level comparison into a
response curve, and would say where on that curve a company's document-handling practice needs to
sit. The protocol from the multilingual parsing benchmark \parencite{mdpbench2026} is a reasonable
template.

\section{Production Hardening}
\label{sec:fw-production}

What a company would additionally need: throughput and batching, audit logging of every extraction
with its provenance, and the review interface that the precision-over-recall error profile
(Section~\ref{sec:results-heldout}) makes viable.

One specific and cheap addition belongs here. The validate-and-repair stage does not currently
check the arithmetic identities the invoice standard guarantees --- line amounts summing to the
line total, taxable base plus tax equalling the grand total, the amount due following from the
grand total and any prepayment (Section~\ref{sec:design-validation}). Those identities are
deterministic, require no model, and would convert a class of silent misreading into a flagged
document. They would not improve any score reported in this thesis, which is precisely why they
were not built; they would improve the reviewability of the output, which is what a company would
actually be buying. The prototype of
Chapter~\ref{ch:implementation} includes a review surface for exactly this reason, but it was
built to support the evaluation, not a company's daily workflow, and the difference between
those two is substantial.
